% Writeup for the Maroon Capital 2026 challenge (a Colonel Blotto variant).
% Build from the repository root: tectonic writeup/writeup.tex
% The figures read the CSV files that scripts/run_pipeline.py writes to results/figure_data/.
\documentclass[10pt,a4paper,twocolumn]{article}

\usepackage[a4paper,top=1.5cm,bottom=1.6cm,left=1.6cm,right=1.6cm,columnsep=0.6cm]{geometry}
\usepackage[lining]{libertinus}
\usepackage{microtype}
\usepackage{amsmath}
\usepackage{booktabs}
\usepackage{enumitem}
\usepackage{titlesec}
\usepackage[font=small,labelfont=bf,skip=3pt]{caption}
\usepackage{xcolor}
\usepackage{pgfplots}
\usepgfplotslibrary{groupplots}
\usetikzlibrary{calc}
\usepackage[hidelinks,pdftitle={Pricing the Battlefields},pdfauthor={Doyoon Kim}]{hyperref}
\pgfplotsset{compat=1.18}

\pagestyle{empty}
\setlist{nosep,leftmargin=1.1em}
\titleformat{\section}{\normalfont\normalsize\bfseries}{\thesection}{0.5em}{}
\titlespacing*{\section}{0pt}{0.8\baselineskip}{0.25\baselineskip}
\hyphenation{FiveThirtyEight Riddler}
\hyphenpenalty=300
\exhyphenpenalty=300
\setlength{\textfloatsep}{8pt plus 2pt minus 2pt}

% Chart colors: the validated categorical slots 1 (blue) and 2 (orange), plus recessive ink and chrome.
\definecolor{ink}{HTML}{0B0B0B}
\definecolor{ink2}{HTML}{52514E}
\definecolor{chrome}{HTML}{C3C2B7}
\definecolor{gridline}{HTML}{E1E0D9}
\definecolor{series1}{HTML}{2A78D6}
\definecolor{series2}{HTML}{EB6834}

\newcommand{\datadir}{../results/figure_data}

\pgfplotsset{
  writeup/.style={
    scale only axis,
    clip marker paths=true,
    font=\sffamily\footnotesize,
    axis lines*=left,
    axis line style={draw=chrome,line width=0.5pt},
    tick style={draw=chrome,line width=0.5pt},
    major tick length=2pt,
    tick label style={text=ink2},
    label style={text=ink2},
    grid style={draw=gridline,line width=0.4pt},
    legend style={draw=none,fill=none,text=ink2,font=\sffamily\footnotesize},
    legend cell align=left,
    table/col sep=comma,
    filter discard warning=false,
  },
  % Keep only the rows whose column #1 equals #2.
  discard if not/.style 2 args={
    x filter/.append code={
      \edef\tempa{\thisrow{#1}}
      \edef\tempb{#2}
      \ifx\tempa\tempb\else\def\pgfmathresult{inf}\fi
    },
  },
}

\begin{document}

\twocolumn[{%
\begin{center}
  {\LARGE\bfseries Pricing the Battlefields}\\[4pt]
  {\normalsize Doyoon Kim\quad\textperiodcentered\quad Maroon Capital 2026 challenge\quad\textperiodcentered\quad dkim927@uchicago.edu}\\[6pt]
  {\large Submission:\enspace\textbf{1, 2, 4, 1, 16, 29, 3, 4, 33, 7}}
\end{center}
\vspace{2pt}
\begin{center}
\begin{tikzpicture}
\begin{axis}[
  writeup,
  width=\linewidth-1.15cm, height=3.7cm,
  xmin=0.5, xmax=10.5, ymin=0, ymax=35,
  xtick={1,...,10}, ytick={0,10,20,30},
  xtick style={draw=none},
  ymajorgrids,
  xlabel={Battlefield}, ylabel={Troops},
  clip=false,
  legend style={at={(0.01,0.96)},anchor=north west},
  reverse legend,
]
  % My allocation: a light bar for the troops spent, and below a line at the price I pay.
  \addplot[ybar,bar width=0.62,fill=series2!14,draw=none,forget plot]
    table[x=battlefield,y=troops] {\datadir/allocation.csv};
  % The chatbot answers, one dot each; equal bids sit side by side.
  \addplot[only marks,mark=*,mark size=1.35pt,mark options={fill=series1,draw=white,line width=0.35pt}]
    table[x expr=\thisrow{battlefield}+\thisrow{offset},y=troops] {\datadir/chatbot_bids.csv};
  % My bid, drawn over the dots so that ties show as dots on the line.
  \addplot[draw=none,mark=none,
    legend image code/.code={\draw[series2,line width=1.5pt] (0cm,0cm) -- (0.36cm,0cm);},
    error bars/.cd,x dir=both,x fixed=0.31,error mark=none,
    error bar style={draw=series2,line width=1.5pt},
  ] table[x=battlefield,y=troops] {\datadir/allocation.csv};
  % Header rows: my troops and the share of the simulated field they beat.
  \addplot[draw=none,mark=none,forget plot,nodes near coords,point meta=explicit symbolic,
    every node near coord/.style={anchor=base,yshift=15pt,font=\sffamily\footnotesize\bfseries,text=ink}]
    table[x=battlefield,y expr=35,meta=troops] {\datadir/allocation.csv};
  \addplot[draw=none,mark=none,forget plot,nodes near coords={\pgfmathprintnumber[fixed,precision=0]{\pgfplotspointmeta}\%},
    point meta=explicit,
    every node near coord/.style={anchor=base,yshift=5pt,font=\sffamily\footnotesize,text=ink2}]
    table[x=battlefield,y expr=35,meta expr=100*\thisrow{share_beaten}] {\datadir/allocation.csv};
  \node[anchor=base east,font=\sffamily\footnotesize,text=ink2] at (axis cs:0.5,35) [yshift=15pt] {My bid};
  \node[anchor=base east,font=\sffamily\footnotesize,text=ink2] at (axis cs:0.5,35) [yshift=5pt] {Beats};
  \legend{Chatbot answer,My bid}
\end{axis}
\end{tikzpicture}
\captionof{figure}{%
My bids (orange line; the light bar shows the troops spent) against 20 chatbot answers (blue dots).
A dot above my line beats me on that battlefield, a dot on it ties, and a dot below it loses.
The header rows give my troops and the share of the whole simulated field that each bid beats.}
\label{fig:bids}
\end{center}
\vspace{4pt}
}]

\noindent\textbf{Summary.}
I target battlefields 2, 3, 5, 6, 9 and 10, three adjacent pairs whose values sum to 35, so winning all six fires every multiplier and scores 80.85.
I outbid all 20 chatbot answers I collected on 6 and 9, match the highest of them on 5, and give up 7 and 8, where most chatbots spend 14 to 31 troops each.
Against a simulated field weighted toward chatbot answers, the allocation beats 93\% of opponents and finishes in the top quartile by wins and by points in all 1,000 mock tournaments.

\begin{figure*}[t]
\centering
\begin{tikzpicture}[baseline=(cand.south)]
\begin{axis}[
  writeup,
  name=cand,
  width=7.0cm, height=3.9cm,
  xmin=0.83, xmax=0.99, ymin=55, ymax=68.5,
  xtick={0.85,0.90,0.95}, xticklabels={85\%,90\%,95\%},
  ytick={56,58,...,68},
  xmajorgrids, ymajorgrids,
  xlabel={Win rate against the center-mix field}, ylabel={Average points per match},
  legend style={at={(0.02,0.04)},anchor=south west},
]
  \addplot[only marks,mark=*,mark size=1.1pt,mark options={fill=chrome,draw=white,line width=0.3pt}]
    table[x=win_rate,y=avg_points] {\datadir/candidates.csv};
  \addplot[series1,line width=1pt,mark=*,mark size=1.3pt,mark options={fill=series1,draw=white,line width=0.3pt},
    discard if not={frontier}{1}]
    table[x=win_rate,y=avg_points] {\datadir/candidates.csv};
  \addplot[only marks,mark=*,mark size=2.6pt,mark options={fill=series2,draw=white,line width=0.8pt},
    discard if not={pick}{1},forget plot,
    nodes near coords={My pick},
    every node near coord/.style={anchor=south west,xshift=2pt,yshift=1pt,font=\sffamily\footnotesize,text=ink}]
    table[x=win_rate,y=avg_points] {\datadir/candidates.csv};
  \legend{Candidate,Pareto frontier}
\end{axis}
\node[anchor=south west,font=\sffamily\bfseries\footnotesize,text=ink] at ([xshift=-1.0cm,yshift=0.1cm]cand.north west) {a};
\end{tikzpicture}%
\hfill
\begin{tikzpicture}[baseline=(tour c1r2.south)]
\begin{groupplot}[
  writeup,
  group style={group name=tour,group size=1 by 2,vertical sep=0.4cm,xticklabels at=edge bottom},
  width=7.0cm, height=1.75cm,
  xmin=0, xmax=0.35, ymin=0,
  xtick={0,0.05,...,0.35}, xticklabels={0\%,5\%,10\%,15\%,20\%,25\%,30\%,35\%},
  ymajorgrids,
]
  \nextgroupplot[ymax=210,ytick={0,100,200}]
    \addplot[ybar,bar width=0.0075,fill=series1,draw=none]
      table[x expr=\thisrow{share_above}+0.005,y=wins] {\datadir/tournaments.csv};
    \draw[ink2,line width=0.6pt] (axis cs:0.25,0) -- (axis cs:0.25,210);
    \node[anchor=north east,font=\sffamily\footnotesize,text=ink2,align=right] at (axis cs:0.246,205) {top-quartile\\cutoff};
    \node[anchor=north east,font=\sffamily\footnotesize,text=ink] at (rel axis cs:1,1) {By wins};
  \nextgroupplot[ymax=1050,ytick={0,400,800},xlabel={Share of the field ranked above my pick}]
    \addplot[ybar,bar width=0.0075,fill=series1,draw=none]
      table[x expr=\thisrow{share_above}+0.005,y=points] {\datadir/tournaments.csv};
    \draw[ink2,line width=0.6pt] (axis cs:0.25,0) -- (axis cs:0.25,1050);
    \node[anchor=north east,font=\sffamily\footnotesize,text=ink] at (rel axis cs:1,1) {By points};
\end{groupplot}
\node[anchor=south,rotate=90,font=\sffamily\footnotesize,text=ink2]
  at ($(tour c1r1.north west)!0.5!(tour c1r2.south west)+(-0.62cm,0)$) {Mock tournaments};
\node[anchor=south west,font=\sffamily\bfseries\footnotesize,text=ink] at ([xshift=-1.0cm,yshift=0.1cm]tour c1r1.north west) {b};
\end{tikzpicture}
\caption{(a) Win rate and average points per match for the 300 candidates.
My pick sits at the knee of the frontier.
(b) My pick's rank in 1,000 mock tournaments, as the share of the field ranked above it.
All 1,000 ranks fall left of the top-quartile cutoff.}
\label{fig:checks}
\end{figure*}

\section{What the scoring rewards}

Each player wins one of $2^{10} = 1{,}024$ sets of battlefields, so I precompute every score and store it as an integer times 100, which makes ties exact.
Four facts shape the game.
\begin{enumerate}
  \item The multipliers compound.
  If every won battlefield has a won neighbor, the base is a multiple of 5, and the set contains 6, the score is $1.5 \times 1.4 \times 1.1 = 2.31$ times the base.
  \item Rule 2 depends on values modulo 5.
  Battlefields 5 and 10 never change it, and the pairs $\{1,4\}$, $\{2,3\}$, $\{6,9\}$ and $\{7,8\}$ cancel.
  \item Winning more can score less.
  In 556 of the 5,120 ways to add one battlefield to a set, the score falls: $\{6,\dots,10\}$ scores 92.4, but $\{1,6,\dots,10\}$ scores 67.1.
  \item You can win a match while losing six battlefields.
  Twenty-six sets of four outscore the other six; $\{2,6,8,9\}$ beats its complement 51.59 to 50.4.
  No set of three or fewer does.
\end{enumerate}
The pairs $\{2,3\}$, $\{5,6\}$ and $\{9,10\}$ use the first two facts at once.
Their values sum to 35, and winning all six scores $35 \times 2.31 = 80.85$, 64\% of the maximum 127.05 from 35 of the board's 55 points.
Battlefields 5 and 10 are the cheap ones to lose, because the base stays a multiple of 5: $\{2,3,6,9,10\}$ still scores 64.68, and $\{2,3,5,6,9\}$ scores 50.82.

\section{The field}

No Blotto allocation wins in isolation, so I model the other applicants.
AI use is allowed, and I expect most applicants to ask a chatbot.
I collected 20 answers by hand, 10 from Claude and 10 from ChatGPT, each from a fresh chat given the problem verbatim.
Most agree on prices (Figure~\ref{fig:bids}): 18 to 27 troops on battlefield 6, 21 to 32 on 9, 14 to 31 on 7 and 8 when they contest them, and at most 3 on 1 and~2.

\begin{table}[t]
\centering\small
\caption{The simulated field.
Range is the share of the field each group may take, and center is its midpoint, rescaled so the shares sum to 100\%.
Win rate is my pick's, against each group with its jittered copies.}
\label{tab:field}
\addfontfeature{Numbers=Tabular}
\begin{tabular}{@{}lrcrr@{}}
\toprule
Group & Entries & Range & Center & Win rate \\
\midrule
Chatbot answers & 20 & 85--98\% & 88.4\% & 98.5\% \\
Riddler entries & 4,510 & 0--5\% & 2.4\% & 73.2\% \\
Rules of thumb & 44 & 0--5\% & 2.4\% & 76.6\% \\
Random splits & 1,000 & 0--2\% & 1.0\% & 67.5\% \\
Counters & 30 & 2--10\% & 5.8\% & 27.7\% \\
\bottomrule
\end{tabular}
\end{table}

I model the field as five groups (Table~\ref{tab:field}).
The group weights are judgment calls, so each group gets a range; the center mix uses the midpoints, and the selection step assumes the worst weights inside the ranges.
Human play comes from 4,510 entries to FiveThirtyEight's Battle for Riddler Nation, which used the same board without the multipliers, so I cap its share at 5\%.
Rules of thumb cover even splits, splits proportional to value, all-in on the top $k$ battlefields, and contiguous blocks whose base is a multiple of 5.
Random splits draw troops from a Dirichlet-multinomial.
Counters stand in for applicants who think a step ahead: six rounds of fictitious play, each adding the five best responses to the pool so far.
Every group except the random one carries a jittered copy of each entry, with 1 to 3 troops moved once or twice and 30\% of the group's weight, so the search cannot exploit one exact number.

\section{The allocation}

Each battlefield has a price list: the share of the field that a given number of troops beats.
I buy where the list jumps and stop one troop past the jump.
\begin{itemize}
  \item \textbf{6 and 9.} 29 and 33 troops top every chatbot answer and beat 97\% and 98\% of the whole field.
  \item \textbf{5.} 16 troops match the highest chatbot bid and beat 91\%.
  \item \textbf{2 and 3.} Six troops in total beat 81\% and 77\%, because 38 of the 40 chatbot bids on them are 3 troops or fewer.
  \item \textbf{10.} 7 troops beat 58\%.
  Losing 10 costs its face value, never the 1.4 multiplier.
  \item \textbf{7 and 8.} 3 and 4 troops concede the two battlefields where the chatbots overpay, yet still beat the 32\% and 20\% of opponents who abandon them.
  \item \textbf{1 and 4.} One troop each ties the many one-troop bids.
\end{itemize}

\begin{table}[t]
\centering\small
\caption{My most common outcomes against the center-mix field.}
\label{tab:outcomes}
\addfontfeature{Numbers=Tabular}
\begin{tabular}{@{}lrrr@{}}
\toprule
Battlefields I win & Base & Score & Matches \\
\midrule
2, 3, 5, 6, 9, 10 & 35 & 80.85 & 20.8\% \\
2, 3, 5, 6, 7, 9, 10 & 42 & 69.30 & 8.8\% \\
2, 3, 5, 6, 9 & 25 & 50.82 & 7.6\% \\
1, 2, 3, 5, 6, 9, 10 & 36 & 59.40 & 5.8\% \\
1, 2, 3, 4, 5, 6, 9 & 30 & 62.37 & 4.6\% \\
1, 2, 5, 6, 8, 9 & 31 & 51.15 & 4.4\% \\
\bottomrule
\end{tabular}
\end{table}

Table~\ref{tab:outcomes} lists my most common outcomes.
The full target set leads at 21\% of matches.
Winning an extra battlefield by accident costs points: adding 7 drops the score to 69.3 and adding 1 drops it to 59.4, because each breaks the multiple of 5.
The troop on battlefield 1 still earns its place, because ties there decide matches.
Against the ChatGPT answer 1, 1, 3, 4, 1, 1, 21, 19, 28, 21, my troop ties theirs and holds their base at 29, so I win 50.82 to 36.5.
Without it they also take battlefield 1, their base reaches 30, their 1.4 multiplier fires, and I lose 50.82 to 52.5.
The same troop also works the other way: against the ChatGPT answer 0, 0, 0, 0, 15, 1, 25, 31, 1, 27, it wins battlefield 1, which pairs with my win on 4 to keep my base at 30, and I win 62.37 to 45.5; without it my base is 29 and I lose 42.9 to 45.5.
Moving that troop to any other battlefield cuts my win rate from 93.0\% to 85.7\% or less.
My won set has a base that is a multiple of 5 in 52\% of matches and contains 6 in 97\%, against 31\% and 46\% for the average chatbot answer.
Head to head, the pick beats all 20 chatbot answers I collected, 10 of 10 from each model.

\section{Search and robust selection}

I rate an allocation $x$ against a weighted field by
\begin{equation*}
  R_\lambda(x) = \lambda\, W(x) + (1-\lambda)\, P(x) / 127.05,
\end{equation*}
where $W$ is the win rate and $P$ the average points per match.
The organizers rank on an unknown mix of wins and points, so $\lambda$ is unknown as well.

I search with simulated annealing over integer allocations, using 32 chains of 3,000 steps that move 1 to 5 troops at a time, then polish with steepest ascent over every shift of 1 to 3 troops until none helps.
I run the search under the center mix and under every extreme mix of the ranges, for $\lambda \in \{0.25, 0.5, 0.75\}$, add two searches that maximize the worst case below, and keep the 300 best distinct candidates.
For speed, each search rates against a stratified sample of 2,000 opponents, and I re-rate the candidates on the full pool of 10,208 entries.
The counters come from the same search, run at $\lambda = 0.5$, keeping the five best results each round that differ by at least 6 troops.

A pessimist then chooses the group weights $w$ inside the ranges $B$ and $\lambda \in [0.25, 0.75]$ to minimize each candidate's rating, and I submit the candidate with the best worst case:
\begin{equation*}
  x^\star = \arg\max_x \; \min_{w \in B,\, \lambda} \; \textstyle\sum_g w_g R_{\lambda,g}(x).
\end{equation*}
The rating is linear in $w$, so a greedy fill gives the pessimist's exact minimum, and the worst $\lambda$ is always an endpoint.
My pick's worst case is 0.598.
Ten candidates sit within 0.005 of it, in three shapes: mine and three close variants, two that concede battlefield 6 and buy 8 and 9, and four that buy 7 and 8 instead of 6 and 9.
All ten finish in the top quartile of every mock tournament I ran on them, so the worst case breaks the tie.
My shape also beats more of the Riddler humans (73 to 74\%, against 60 to 65\% for the others) and scores about three more points per match than the shape that concedes 6.
Figure~\ref{fig:checks}a shows the trade-off the pessimist forces: the candidate with the most wins beats 98.5\% of the field but scores 55.7 points per match, while my pick gives up 5.6 points of win rate to score 65.2.

I considered two alternatives and rejected both.
An equilibrium strategy mixes over many allocations, but I submit one allocation into a field that does not play the equilibrium.
Reinforcement learning would chase the same objective with more noise: the game is a single move, and I can score any allocation against the whole field in milliseconds.


\section{Validation}

\textbf{Reproducibility.}
One command runs the whole pipeline in about 90 seconds with fixed random seeds, and a second script reruns the weight comparison below, so every number here reproduces exactly.
Seventy-four tests check it, including a cross-check of all 1,024 scores against an independent exact-fraction implementation and a check that the greedy worst case matches brute force over every extreme mix.

\textbf{Mock tournaments.}
I drew 1,000 fields of 100 to 300 applicants, with group weights uniform inside the ranges, and played full round robins.
My pick finished in the top quartile by wins and by points in all 1,000 (Figure~\ref{fig:checks}b).
The median share of the field ranked above it is 3.5\% by wins, and in 72\% of tournaments no applicant scored more total points.

\textbf{Held-out data.}
A rerun that never saw ChatGPT's answers or Riddler round 4 chose 2, 3, 3, 6, 14, 28, 3, 4, 31, 6, with the same shape: heavy on 5, 6 and 9 and light elsewhere.
It beats 80\% of the unseen ChatGPT answers, against 100\% of the Claude answers it trained on, and its points per match fall from 76.9 to 58.6.
On Riddler round 4 its win rate holds at 73\%, against 74\% on the rounds it saw.

\textbf{Counters.}
My exact allocation has a counter: 0, 4, 3, 4, 22, 0, 12, 10, 36, 9 gives up 1 and 6 and outbids me almost everywhere else, and my pick wins none of 51 matches against it and 50 jittered copies of it.
No entry in the simulated field lies within 10 troops of that counter, so beating me this way takes an applicant who targets my numbers.

\textbf{Sensitivity to the weights.}
My first run used wider ranges: chatbots 70 to 95\%, Riddler up to 20\%, rules of thumb up to 15\%, and counters 5 to 15\%.
It chose 1, 4, 6, 1, 7, 28, 3, 5, 33, 12, the same three pairs with more troops on 10 and fewer on 5.
I narrowed the ranges because Riddler Nation had no multipliers, so its players are a weak model for this game.
Under the old ranges, with the same simulated pool, the new pick still finishes in the top quartile of all 1,000 mock tournaments, with a median of 5.7\% of the field above it by wins against 10.3\% for the old pick.
The cost is robustness to strategic players: the old pick beats 57\% of the counters, while the new pick beats 28\%.

\section{Assumptions and risks}

I read the rules this way: battlefields sit on a line, so 1 and 10 are not adjacent; equal totals are a win for neither player; and total points means the sum of my own scores.
The chatbot share is my biggest judgment call, and the pick leans on it: it beats 98.5\% of chatbot answers but only 28\% of the counters, so a field full of applicants who run their own optimization would hurt it.
Against the Riddler humans it still wins 73\%.
The chatbot sample is small: 20 answers from two models and one prompt, and real applicants may edit what a chatbot gives them.
The mock tournaments draw from the pool that produced the pick, so they flatter it; the held-out run, where the win rate against unseen ChatGPT answers fell to 80\%, gives the fairer estimate.

\end{document}
